\section{Background}\label{sec:background}

The framework built on three established models, each contributing a distinct component of the synthetic option generation pipeline: the JumpHMM produced realistic multi-asset price paths, the Heston stochastic variance process provided a mechanism for converting price dynamics into implied volatility, and the CRR binomial tree translated IV into American option prices with early exercise.

\subsection{Jump Hidden Markov Model}

The Jump Hidden Markov Model (JumpHMM)~\cite{jumphmm2025} was developed for generating synthetic equity time series that simultaneously preserved three stylized facts of financial returns: heavy tails, negligible linear autocorrelation, and persistent volatility clustering~\cite{cont2001}. The model discretized continuous excess growth rates into $N$ states via equal-probability quantile bins of a fitted Laplace distribution. The hidden state $S_t \in \{1, \ldots, N\}$ at each timestep governed the distribution of the excess growth rate $G_t$; each state $k$ emitted observations from a location-scale Student-$t$ distribution: $G_t \mid S_t\!=\!k \;\sim\; \mu_k + \sigma_k \cdot t_\nu$, where $\mu_k$ and $\sigma_k$ were learned per-state location and scale parameters and $\nu$ was the shared degrees-of-freedom parameter. At each timestep, with probability $\epsilon$ a Poisson$(\lambda)$-distributed jump forced the chain into tail states (bottom or top $N_{\text{tail}}$ states), creating sustained dwell periods in extreme-volatility regimes that produced the volatility clustering characteristic of empirical return series. For multi-asset generation, cross-asset dependence was imposed via a Student-$t$ copula~\cite{demarta2005} that preserved each asset's marginal distribution while introducing realistic tail dependence; this copula structure ensured that joint tail events, in which multiple assets simultaneously occupied extreme states, occurred more frequently than a Gaussian copula would imply.

\subsection{Heston stochastic variance model}

The Heston model~\cite{heston1993} specified the instantaneous variance $v$ of the underlying asset as a mean-reverting square-root process:
\begin{equation}\label{eq:heston_standard}
    dv = \kappa(\theta - v)\,dt + \sigma_v \sqrt{v}\,dW_v
\end{equation}
where $\kappa$ was the mean-reversion speed controlling how quickly the variance returned to its long-run level, $\theta$ was the long-run variance target, and $\sigma_v$ was the vol-of-vol governing the amplitude of stochastic fluctuations in $v$. In the standard formulation, the asset price and variance processes were coupled through a correlation $\rho$ between their driving Brownian motions, producing the leverage effect in which negative returns coincided with rising volatility. The key property for our purposes was mean-reversion: the variance process was pulled toward $\theta$, so by controlling $\theta$ we could control the equilibrium level of implied volatility. In the standard model $\theta$ was a constant; our extension made $\theta$ time-varying and state-dependent, linking it to the JumpHMM regime state, contract characteristics, and market conditions (Section~\ref{sec:theta_function}).

\subsection{CRR binomial tree}

The CRR model~\cite{cox1979} priced options on a recombining binomial lattice with up/down factors $u = e^{\sigma\sqrt{\Delta t}}$, $d = 1/u$ and risk-neutral probability $p = (e^{r\Delta t} - d)/(u - d)$. American options were priced by backward induction with an early-exercise check at each node, making the CRR tree a standard tool for contracts where closed-form solutions do not exist. The CRR tree was a natural choice when IV was provided exogenously, as it required only a single volatility input per contract and handled early exercise without the need for Monte Carlo simulation or PDE solvers. Convergence to the continuous-time Black-Scholes price was guaranteed as the number of tree steps increased, and the computational cost scaled linearly in the number of steps, making the tree efficient for pricing large cross-sections of contracts across multiple strikes and expirations.
